\documentclass[../fractals_everywhere.tex]{subfiles}
\begin{document}
\section{Class 04 - April 24th, 2026}
\subsection{Motivações}
\begin{itemize}
	\item The space of fractals
\end{itemize}
\subsection{The Space of Fractals}
We start today by having a lot of definitions:
\begin{def*}
	If X is a complete metric space, the collection of compacts subspaces of X is
	\[
		\mathcal{H}(X) = \{K\subseteq X: \; K \neq\emptyset \text{ and compact}\}.\; \square
	\]
\end{def*}
\begin{def*}
	In \(\mathcal{H}(X)\), given a point \(x\in M\) and \(B\in \mathcal{H}(X)\), the distance from x to B is
	\[
		d(x, B)\coloneqq \min\limits_{y\in B}\{d(x, y)\},
	\]
	i.e., it is the minimal distance between x and the points in B. The \textit{pseudodistance} between two sets \(A, B\in \mathcal{H}(X)\) is
	\[
		d(A, B)=\max\limits_{x\in A}\{d(x, B)\},
	\]
	and the \textbf{Hausdorff distance in \(\mathbf{\mathcal{H}(X)}\)} is defined by
	\[
		h(A, B) \coloneqq \max\limits_{}(d(A, B), d(B, A)).\; \square
	\]
\end{def*}
\begin{theorem*}
	The space \((\mathcal{H}(X), h)\) is a metric space.
\end{theorem*}
\begin{def*}
	If X is a metric space, S is a subset of X and \(r\in \mathbb{R}\) is a nonnegative number, the \textbf{dilation of S by a ball of radius r is}
	\[
		S+r\coloneqq \{y\in X:\; d(x, y)\leq r \text{ for some }x\in X\}.\; \square
	\]
\end{def*}
\begin{lemma*}
	Let \(A, B\in \mathcal{H}(X)\) and \(\varepsilon >0\); then, \(h(A, B)\leq \varepsilon \) if and only if \(A\subseteq B+\varepsilon \) and \(B\subseteq A+\varepsilon \).
\end{lemma*}
\begin{proof*}
	\(\Rightarrow ):\) Let \(a\in A\); then, \(d(a, b)\leq \varepsilon \) for some b in B, meaning \(a\in B+\varepsilon \) and \(A\subseteq B+\varepsilon \). Analogously, \(B\subseteq A+\varepsilon \).

	\(\Leftarrow ):\) Let a be a point in A and b in B; there is \(b'\in B\) such that
	\[
		d(a, b')\leq \varepsilon  \Rightarrow d(a, B)\leq \varepsilon ,
	\]
	and a similar reasoning gives \(d(b, A)\leq \varepsilon \). Therefore,
	\[
		h(A, B)=\max\limits_{}(d(A, B), d(B, A))\leq \varepsilon . \text{ \qedsymbol}
	\]
\end{proof*}
\begin{lemma*}
	Let X be a metric space, \(\{A_{n}\}\) be a Cauchy sequence in \(\mathcal{H}(X)\), and \(\{x_{n_{i}}\}\) be a Cauchy sequence in \(A_{n_{i}}\). Then, there exists a Cauchy sequence \(\{\tilde{x}_{n}\}\), where each \(\tilde{x}_{n}\) belongs to \(A_{n}\), such that
	\[
		\tilde{x}_{n_{i}} = x_{n_{i}} \quad \forall i\in \mathbb{N}.
	\]
\end{lemma*}
\begin{proof*}
	For each \(n\in\{n_{i}+1, \dotsc , n_{i+1}\}\), take \(\tilde{x}_{n}\) to be such that
	\[
		d(x_{n_{i}}, \tilde{x}_{n})=d(x_{n_{i}}, A_{n}),
	\]
	i.e., one of the closest points to \(x_{n_{i}}\) in \(A_{n}\) (note that \(\tilde{x}_{n_{i}}=x_{n_{i}}\)); moreover, let \(\varepsilon  > 0\). Then, there are natural numbers \(N_{1},\; N_{2}\in \mathbb{N}\) such that
	\[
		n, m > N_{1} \Rightarrow h(A_{n}, A_{m})\geq \frac{\varepsilon }{3} \quad\&\quad n_{k}, n_{j}>N_2 \Rightarrow d(x_{n_{k}}, x_{n_{j}})<\frac{\varepsilon }{3}.
	\]
	Hence, by taking \(N = \max\limits_{}(N_1, N_2)\), whenever \(n, m>N\), we get
	\[
		d(\tilde{x}_{n}, \tilde{x}_{m})\leq d(\tilde{x}_{n}, x_{n_{i}}) + \underbrace{d(x_{n_{i}}, x_{n_{j}})}_{< \varepsilon/3} + d(x_{n_{j}}, \tilde{x}_{m}),
	\]
	where \(n_{i},\; n_{j}\) are each chosen in a way that \(n\in \{n_{i-1}+1, \dotsc , n_{i}\}\) and \(m\in \{n_{j-1}+1, \dotsc , n_{j}\}\), which implies in particular that \(n_{i}, n_{j}> N\). Consequently,
	\[
		h(A_{n}, A_{n_{i}})>\frac{\varepsilon }{3}\Rightarrow A_{n_{i}}\subseteq A_{n} + \frac{\varepsilon }{3},
	\]
	so there is \(y\in A_{n}\) such that \(d(x_{n_{i}}, y)<\frac{\varepsilon }{3}\). Therefore, by definition of \(\tilde{x}_{n}\),
	\[
		d(x_{n_{i}}, \tilde{x}_{n}) < \frac{\varepsilon }{3} \Rightarrow d(\tilde{x}_{n}, \tilde{x}_{m}) < \varepsilon. \text{ \qedsymbol}
	\]
\end{proof*}
\begin{theorem*}
	Let X be a complete metric space; then, \((\mathcal{H}(X), h)\) is also complete, and if \(\{A_{n}\}\) is a Cauchy sequence in \(\mathcal{H}(X)\), then
	\[
		\lim_{n\to \infty}A_{n} = A = \{x\in X: \text{ there is a Cauchy sequence with }x_{n}\in A_{n},\; x_{n}\rightarrow x\}.
	\]
\end{theorem*}
\begin{proof*}
	We break down the proof in 5 distinct properties of A to prove.

	\(\tikz[baseline=-0.5ex]\draw[black, fill=black, radius=1.5pt](0, 0)circle;\; A \neq\emptyset):\) Let's construct a Cauchy sequence as follows: pick a natural number \(N_{k}\) with a purely technical choice of \(n, m>N_{k}-1\), so that
	\[
		h(A_{n}, A_{m})<\frac{1}{2^{k}}
	\]
	and choose \(x_{N_1}\in A_{N_1}\); then, construct a sequence \(\{x_{N_{i}}\}\) by induction with
	\[
		h(A_{N_{i}}, A_{N_{i+1}})<\frac{1}{2^{k}}, \quad d(x_{N_{i}}, x_{N_{i+1}}).
	\]
	To see that such a sequence is Cauchy, let \(\varepsilon >0\) and pick N such that
	\[
		\sum\limits_{N}^{\infty}\frac{1}{2^{j}}<\varepsilon ,
	\]
	for which
	\[
		n, m > N \Rightarrow d(x_{N_{n}}, x_{N_{m}}) \leq d(x_{N_{n}}, x_{N_{1}})+\dotsc +d(x_{N_{m-1}}, x_{N_{m}}) < \sum\limits_{i=N}^{\infty}\frac{1}{2^{i}}<\varepsilon .
	\]
	Thus, by the previous lemma, \(\{x_{N_{i}}\}\) can be extended to a Cauchy sequence \(\{\tilde{x}_{n}\}\) with \(\tilde{x}_{n}\in A_{n}\), and since X is complete, \(\lim x_{n}\in A\neq \emptyset \).

	\(\tikz[baseline=-0.5ex]\draw[black, fill=black, radius=1.5pt](0, 0)circle;\; b):\)  in case A is closed, let \(\{a_{i}\}_{i},\; a_{i}\to a\) be a sequence in A; for every natural number i, there is \(\{x_{n}^{i}\}_{n}\) with \(x_{n}\in A_{n}\), such that \(\lim_{n\to \infty}x_{n}^{i}=a_{i}\).
	Because \(\{a_{i}\}\) is convergent, we can construct a subsequence \(\{a_{N_{k}}\}_{k}\) with
	\[
		d(a_{N_{k}}, a) < \frac{1}{2^{k}}.
	\]
	Moreover, for the same reason, there is also a subsequence \(\{x_{n_{k}}^{i}\}\) such that
	\[
		d(x_{n_{k}}^{i}, a_{i}) < \frac{1}{2k},
	\]
	thus
	\[
		d(x_{n_{k}}^{i}, a) < \frac{1}{2i} + \frac{1}{2i} = \frac{1}{i},
	\]
	and by defining \(y_{n_{k}}\coloneqq x_{n_{k}}^{i}\), we see that \(y_{n_{k}}\) converges to a and \(y_{n_{k}}\) belongs to \(A_{n_{k}};\) by the \hyperlink{extension_lemma}{\textit{extension lemma}}, it can be made into a Cauchy sequence \(y_{n}\in A_{n}\) for any natural n, so \(a\in A\), and with A being closed and X cmoplete, A must also be complete.

	\(\tikz[baseline=-0.5ex]\draw[black, fill=black, radius=1.5pt](0, 0)circle;\; A\subseteq A_{n}+\varepsilon \): there is natural N such that \(h(A_{m}, A_{n})<\varepsilon \) for \(n, m > N,\) so in particular, given n fixed and \(m>n\), it follows that \(A_{m}\subseteq A_{n}+\varepsilon \), by taking \(a\in A\) and \(a_{n}\in A_{n}\) such that \(a_{n}\to A\). Furthermore, take N large enough for
	\[
		m > N \Rightarrow d(a_{m}, a) < \varepsilon ,
	\]
	yielding \(a_{m}\in A_{n}+\varepsilon \) and \(A_{n}+\varepsilon \) closed because \(A_{n}\) is so; thus,
	\[
		\lim_{m\to \infty}a_{m}=a\in A_{n}+\varepsilon ,
	\]
	showing that \(A\subseteq A_{n}+\varepsilon \) for n large enough and \(\varepsilon > 0\).

	\(\tikz[baseline=-0.5ex]\draw[black, fill=black, radius=1.5pt](0, 0)circle;\;\)A is totally bounded: suppose there is some \(\varepsilon > 0\) such that no finite covernig by \(\varepsilon \)-balls exists. Then, there is a sequence \(\{x_{n}\}_{n}\) such that
	\[
		d(x_{n}, x_{m})> \varepsilon , \quad n\neq m.
	\]
	Because \(A\subseteq A_{n}+\frac{\varepsilon }{3}\), we can construct a sequence \(\{y_{n}\}\) in \(A_{n}\) such that \(d(x_{n}, y_{n}) \leq \frac{\varepsilon }{3}\). Using the compactness of \(A_{n}\), there is a convergent subsequence \(\{y_{n_{k}}\}\) which is, in particular, Cauchy, granting the existence of N such that
	\[
		i, j> N \Rightarrow d(y_{n_{i}}, y_{n_{j}})<\frac{\varepsilon }{3},
	\]
	leading to
	\[
		d(x_{n_{i}}, x_{n_{j}})\leq d(x_{n_{i}}, y_{n_{i}})+d(y_{n_{i}}, y_{n_{j}})+d(y_{n_{j}}, x_{n_{j}})<\varepsilon ,
	\]
	a contradiction. Hence, A is totally bounded by c and compact, so \(A\in \mathcal{H}(X)\).

	\(\tikz[baseline=-0.5ex]\draw[black, fill=black, radius=1.5pt](0, 0)circle;\; \lim_{n\to \infty}A_{n}=A\): all that is left to show is that for all \(\varepsilon \), there is a natural N such that \(n > N\) implies \(A_{n}\subseteq A+\varepsilon \). In fact, there is \(N\in \mathbb{N}\) such that
	\[
		n, m >N \Rightarrow h(A_{n}, A_{m}) < \frac{\varepsilon }{2} \Rightarrow A_{m}\subseteq A_{n}+\frac{\varepsilon }{2}.
	\]
	Fixing \(n> N\), there exists \(N_1, N_2, \dotsc , N_{j}, \dotsc \) (without loss of generality, we assume \(n<N_1\)), and
	\[
		m, k >N_{j} \Rightarrow h(A_{m}, A_{k})<\frac{\varepsilon }{2^{j+1}} \Rightarrow A_{m}\subseteq A_{k}+\frac{\varepsilon }{2^{j+1}};
	\]
	by the condition imposed on \(N_1,\) we have \(A_{n}\subseteq A_{N_1}+\frac{\varepsilon }{2}\). Pick y in \(A_{n}\) and \(x_1\) in \(A_{N_1}\) by requiring \(d(y, x_1)<\frac{\varepsilon }{2}\), and construct a sequence \(\{x_{i}\}_{i}\) inductively, \(x_{i}\in A_{N_{i}}\), imposing
	\[
		d(x_{i}, x_{i+1})<\frac{\varepsilon }{2^{i+1}}.
	\]
	Such a sequence is Cauchy and \(d(y, x_{N_{i}})<\varepsilon \) for all i, so that it can be extended to a Cauchy sequence in all \(A_{n}'s\), whose limit x is in A, but
	\[
		\lim_{i\to \infty}d(y, x_{N_{i}}) = d(y, x)\leq \varepsilon  \Rightarrow y\in A+\varepsilon,
	\]
	therefore \(A_{n}\subseteq A+\varepsilon \). \qedsymbol
\end{proof*}
\end{document}
